\chapter{Nexo}

O aplicativo proposto por esse trabalho, denominado ``Nexo'', se propõe a ser uma
aplicação para dispositivos móveis, de forma a preencher essas lacunas
demonstradas tanto pelo questionário, quanto também tentando preencher as
questões analisadas em outros sistemas similares. Como mencionado
anteriormente, o sistema é direcionado para pessoas com TAB, embora, ele também
possa atuar isoladamente como um aplicativo de registro de humor, e portanto, o
nome reflete essa escolha. A lógica é, de um aplicativo que consegue conectar
diversos pontos, que em um primeiro momento parecem soltos: sono, rotina,
práticas, humor, gatilhos.

A partir das demandas identificadas na pesquisa e dos princípios estabelecidos
para este projeto, delineia-se um sistema de monitoramento de humor centrado no
usuário, com ênfase em personalização, privacidade e utilidade reflexiva. O
funcionamento do sistema parte da premissa de que o humor não é uma variável
universal, mas sim composto por elementos subjetivos que variam entre
indivíduos. Assim, fundamentalmente para cada entrada no diário, o usuário pode
indicar que componentes de humor ele percebe em vigência naquele momento, como
``raiva'', ``tristeza'', ``felicidade'', etc., elementos que ele considera mais
representativos de seu estado emocional, permitindo um mapeamento mais fiel da
sua experiência.

O monitoramento, pode ser gerado em diversos momentos do dia, o que torna o
aplicativo mais flexível, ao invés de uma única entrada no dia. O usuário pode
também anotar para cada momento, algo relevante, através de escrita livre, e
usar uma escala que possa indicar o que ele percebe como indicadores dos níveis
de energia, estresse e ansiedade. Associado a isso, o sistema permitirá a
entrada de gatilhos e eventos significativos, sejam positivos ou negativos, que
poderão ser posteriormente correlacionados com os estados de humor, permitindo
identificar padrões ou ciclos. Esses dados podem ser correlacionados/vinculados
com ``práticas'' ou ``intervenções'', como fármacos, psicoterapia, atividades
físicas.

Outro ponto essencial é a exportação de dados em formatos acessíveis, como o
formato CSV \cite{Shafranovich_2005} \footnote{Formato de arquivo onde cada
coluna é separado por um delimitador, por exemplo: vírgulas, ponto-e-vírgula,
apresentado no RFC 4180}, garantindo que o usuário tenha controle e mobilidade
sobre suas informações, algo frequentemente negligenciado em sistemas fechados.
Além disso, a ferramenta é por inteiro \textit{software} livre, garantindo
transparência de funcionamento, possibilidade de auditoria e fomento a
melhorias comunitárias, especialmente importantes no contexto de tecnologias
sensíveis à saúde mental. Essas características formam a base de um sistema
comprometido com o bem-estar, o respeito à individualidade e a soberania dos
dados pessoais.

Com base nas funcionalidades listadas, pode-se definir os requisitos funcionais
e não funcionais no qual o sistema deve necessitar.

\section{Classificação dos Requisitos}

Como observado por \cite[p.~58]{sommerville:2010}, ``Os requisitos de um sistema são
as descrições dos serviços que o sistema deve prestar e as restrições à sua
operação''. Portanto, requisitos são declarações articuladas descrevendo de
forma clara o que o sistema deve ser capaz de fazer para satisfazer as
necessidades postas anteriormente. São definidas em duas categorias: requisitos
funcionais que descrevem o comportamento e funções do sistema e requisitos não
funcionais que descrevem os critérios específicos que podem ser usados para
avaliar o bom funcionamento do sistema. A separação entre requisitos funcionais
e não funcionais permite uma estrutura clara para priorização, testes e
validação do projeto, respeitando tanto os objetivos técnicos quanto éticos da
proposta. Logo a seguir, serão abordados com mais detalhes cada um deles.

\subsection{Requisitos funcionais}

Os requisitos funcionais dizem respeito às funcionalidades centrais do sistema,
ou seja, declarações de tudo o que o sistema deve e está proposto a fazer e
reagir. Sommerville reforça esse aspecto:

\begin{quote}

  Os requisitos funcionais de um sistema descrevem o que ele deve fazer. Eles
  dependem do tipo de software a ser desenvolvido, de quem são seus possíveis
  usuários e da abordagem geral adotada pela organização ao escrever os
  requisitos. Quando expressos como requisitos de usuário, os requisitos
  funcionais são normalmente descritos de forma abstrata, para serem
  compreendidos pelos usuários do sistema.

  \cite[p.~59]{sommerville:2010}
\end{quote}

Portanto, diante das funcionalidades previstas, que foram brevemente listadas,
no Quadro~\ref{tab:requisitos_funcionais} é apresentado os requisitos que o
sistema deve fornecer:

\newpage

% TODO: Add this to the list of tables
\begin{longtable}{| p{.20\textwidth} | p{.30\textwidth} | p{.40\textwidth} |}
    \caption{Requisitos funcionais}
    \label{tab:requisitos_funcionais} \\
        \hline
        N° & Requisito & Descrição \\
        \hline
        \endfirsthead
        \hline
        N° & Requisito & Descrição \\
        \hline
        \endhead
        \endfoot

        [RF01] &
        Login &
        O sistema deve ser capaz de realizar login com email e senha \\

        \hline

        [RF02] &
        Redefinir senha &

        O sistema permite que o usuário redefina a sua senha, mediante informar
        o email usado \\

        \hline

        [RF03] &
        Cadastro &
        O sistema deve permitir que o usuário crie uma conta \\

        \hline

        [RF04] &
        Desconectar &
        O sistema permite que o usuário faça logout do sistema \\

        \hline

        [RF05] &
        Manter humor &
        O usuário deve poder registrar e atualizar seu estado emocional
        conforme a necessidade \\

        \hline

        [RF06] &
        Registrar gatilho &
        O sistema deve permitir associar sentimentos ou eventos específicos que
        influenciam o humor \\

        \hline

        [RF07] &
        Visualizar histórico &
        O usuário deve ter acesso aos registros anteriores de humor e demais
        tipos de dados como Gatilho, Sono, etc. \\

        \hline

        [RF08] &
        Visualizar detalhes do dia &
        O sistema deve permitir expandir e consultar dados específicos de um
        determinado dia \\

        \hline

        [RF09] &
        Visualizar registros de humor recente &
        O usuário deve poder visualizar os ultimos registros de humor como
        forma de ler \\

        \hline

        [RF10] &
        Manter registro de sono &
        O sistema deve permitir que o usuário possa cadastrar registros de sono \\

        \hline

        [RF11] &
        Manter ações de cuidado &
        O usuário deve poder cadastrar ações de cuidado (atividade, medicamentos, consultas) \\

        \hline

        [RF12] &
        Exportar dados &
        O sistema deve oferecer uma funcionalidade para exportar os dados
        registrados \\

        \hline

        [RF13] &
        Registrar que um medicamento foi tomado &
        O usuário pode marcar que um determinado medicamento foi tomado
        mediante que foi cadastrado \\

        \hline

        [RF14] &
        Associar humor à um gatilho &
        O sistema deve permitir que o usuário vincule um humor específico com
        um gatilho \\

        \hline

        [RF15] &
        Associar consulta com um humor ou gatilho &
        O sistema deve permitir vincular uma consulta específica a um
        gatilho ou humor registrado \\

        \hline

        [RF16] &
        Gerar um relatório com os últimos dados &
        O usuário deve selecionar um período e gerar/exportar um relatório com
        os últimos dados  \\

        \hline
        \multicolumn{3}{c}{\small Fonte: Elaborado pelo autor (2026)} \\
\end{longtable}

\subsection{Requirementos não funcionais}

De acordo com \cite[p.~89]{sommerville:2010}: ``Requisitos não funcionais são
restrições sobre os serviços ou funções oferecidas pelo sistema''. Ou seja,
definem os critérios de qualidade que o sistema deve seguir, abrangendo
aspectos como segurança, usabilidade, desempenho e manutenção. Neste projeto,
destacam-se especialmente a preservação da privacidade do usuário, a
transparência proporcionada pelo código aberto e a possibilidade de utilização
offline, o que garante acessibilidade mesmo em contextos de conectividade
limitada.

\begin{longtable}{| p{.20\textwidth} | p{.30\textwidth} | p{.40\textwidth} |}
    \caption{Requisitos não funcionais}
    \label{tab:requisitos_nao_funcionais} \\
        \hline
        N° & Requisito & Descrição \\
        \hline
        \endfirsthead
        \hline
        N° & Requisito & Descrição \\
        \hline
        \endhead
        \endfoot

        [RNF01] &
        Usabilidade &

        O sistema deve possuir uma interface intuitiva e acessível para
        usuários leigos em tecnologia, com navegação fluída e confirmações
        visuais após cada interação importante \\

        \hline

        [RNF02] &
        Desempenho &

        O tempo de resposta entre a entrada do usuário e a resposta do sistema
        deve ser inferior a 1 segundo; O carregamento do histórico semanal e
        visualizações de ciclos anteriores deve ocorrer em no máximo 2
        segundos. \\

        \hline

        [RNF03] &
        Segurança &

        Os dados do usuário devem ser armazenados de forma criptografada; O
        sistema deve exigir autenticação segura para acesso; A exportação de
        dados deve sempre exigir a confirmação do usuário \\
        \hline

        [RNF04] &
        Confiabilidade &

        O sistema deve garantir a persistência dos dados, mesmo em casos de
        falha de energia ou conexão; Os dados registrados devem ser armazenados
        automáticamente e sincronizados quando ouvir conexão \\

        \hline

        [RNF05] &
        Portabilidade &

        Todos os dados do usuário podem ser exportados a qualquer momento, em
        um formato de texto que permita que ele possa ler e interpretar,
        mediante compreensão mínima da linguagem/formato utilizado \\

        \hline

        [RNF06] &
        Privacidade &

        Nenhuma informação sensível do usuário deve ser compartilhada com
        terceiros e o sistema deve fornecer uma política de privacidade clara e
        acessível; Quaisquer alterações à mesma demanda que o usuário seja
        notificado \\

        \hline
        \multicolumn{3}{c}{\small Fonte: Elaborado pelo autor (2026)} \\
\end{longtable}

Com essas definições em mãos, pode-se avançar para a etapa de modelagem do
sistema, usando os requerimentos apresentados como fundação para o sistema.

\section{Modelagem}

De acordo com \cite{guedes:2022}, a modelagem de sistemas é uma etapa fundamental do
processo de desenvolvimento de software, pois permite representar graficamente
os principais elementos do sistema, seus comportamentos e interações. Essa
representação facilita a comunicação entre desenvolvedores, usuários e demais
envolvidos, servindo tanto como guia para implementação quanto como instrumento
de validação dos requisitos levantados. Através da modelagem, é possível
antever a estrutura do sistema, simular sua dinâmica e assegurar que ele
atenderá às necessidades específicas dos usuários.

Neste subcapítulo, são apresentados dois tipos principais de diagramas da
\textit{Unified Modeling Language} (UML) empregados na modelagem do sistema
proposto: o diagrama de casos de uso, que descreve as funcionalidades esperadas
a partir da perspectiva do usuário; o diagrama de classes, que estrutura os
elementos do domínio em termos de atributos e relacionamentos. Esses diagramas
complementam-se mutuamente, proporcionando uma visão abrangente do
comportamento e da arquitetura do sistema.

\subsection{Casos de Uso}

O diagrama de casos de uso é o mais geral, informal e simples de compreender
\cite{guedes:2022}. A base dele é o levantamento de requisitos e a pesquisa
realizada, e serve como o pontapé para todo o processo de modelagem.

O diagrama de casos de uso, apresenta de forma simples, os atores, que são
representações das entidades que interagem com o sistema (sejam eles usuários
finais ou outros sistemas), e os casos de uso, que são as funcionalidades que o
sistema oferece. Ambos são conectados por linhas de associação. O sistema
também é delimitado por um retângulo, que contém os casos de uso
\cite{ibm:2021}.

Ao final desse processo, essa modelagem é transformada em documentação, e no
apêndice B deste documento, contém a mesma, oferecendo uma visão mais detalhada
para cada caso de uso, seus caminhos alternativos e caminhos de exceções.
Pode-se verificar que esse acaba por ser o diagrama mais basilar do projeto
inteiro \cite{guedes:2022}.

\begin{figure}[h]
\centering
\caption{Diagrama de Casos de Uso}
\label{fig:use-cases}
\includegraphics[width=0.9\linewidth]{use-cases.png}
\source{Elaborado pelo autor (2026)}
\end{figure}

Na Figura~\ref{fig:use-cases}, pode-se observar que apenas um ator interage
com o sistema, o próprio usuário final. Também visualiza-se que cada elipse,
se resume a um requisito funcional, e os casos de uso estão interligados por
setas nomeadas com a palavra \textit{``<<extend>>''} que indica conexões onde a
execução de um caso de uso adjacente pode ocorrer dado uma determinada
condição em adição ao caso de uso inicial \cite{ibm:2021}.

Permeando todas as interações no sistema, o caso de uso mais central, o
``Login'', cujo qual estabelece uma base de autenticação, de forma que garante
que apenas usuários autenticados possam interagir com o sistema.

No caso de uso ``Manter Humor'' representa a funcionalidade central do sistema,
permitindo que o usuário registre seu estado emocional ao longo do tempo. Essa
ação ocorre por meio da inserção direta de uma informação de humor, cujo qual o
usuário pode selecionar entre cinco humores ``basais'', além de elementos que
ajudem a compreender, como escalas indicativas dos respectivos níveis de
estresse, ansiedade e quais componentes do humor o usuário identifica com mais
presença. A manutenção do humor no sistema contempla tanto a criação quanto a
edição e a exclusão desses registros.

O caso de uso ``Registrar Intervenção'' contempla o cadastro de estratégias,
práticas ou sugestões que o usuário utiliza ou deseja utilizar para lidar com
oscilações emocionais. As ações de cuidados, podem incluir desde ações
práticas, como caminhar, até técnicas cognitivas ou expressivas, como
atividades criativas, medicamentos utilizados ou atendimentos com profissionais
da área. Essa funcionalidade contribui para que o sistema sirva como uma
central e colabora para construir um registro mais criterioso, visto que o
usuário reconstrói o que aconteceu naquele determinado ponto no tempo.

O caso de uso ``Registrar Gatilho'' visa permitir que o usuário registre
eventos, sejam eles pensamentos, situações, ou estímulos que tenham
influenciado significativamente seu estado emocional. O objetivo é construir um
histórico de causas associadas a variações no humor, de modo a fomentar a
identificação de padrões recorrentes ou específicos. Essa funcionalidade é
especialmente útil para usuários em contextos terapêuticos ou de
autoconhecimento, favorecendo a análise de relações entre circunstâncias
externas e reações internas.

O caso de ``Registro de sono'', permite o usuário, de forma simples, criar no
sistema um registro da quantidade de horas de sono, e adicionar alguma anotação
breve. Esse elemento corrobora com a construção de indicadores, tanto da média
de sono, quanto também de estimular uma correlação entre sono e os níveis de
energia informado em cada humor.

Por fim, o caso de uso ``Vincular humor a um gatilho'' permite ao usuário
estabelecer vínculos entre os gatilhos identificados e os registros de humor,
formando um mapeamento de ação e reação, ou causa e reação. Essa associação
pode servir tanto para consulta futura quanto para, em um futuro, gerar
indicadores. Ao conectar eventos e suas respostas, a funcionalidade fortalece o
papel do sistema como uma ferramenta alinhada com o EMA, gerando assim uma
narrativa emocional diária.

\subsection{Diagrama de Classes}

O diagrama de classes é um dos principais artefatos da UML, responsável por
evidenciar a estrutura estática do sistema, seus atributos, operações e
relacionamentos. Conforme \cite[p. 101]{guedes:2022}: ``O principal enfoque é
permitir a visualização das classes que o sistema e seus respectivos métodos e
atributos''. Dessa forma, esse diagrama constitui a base da modelagem
estrutural que posteriormente será implementada em código.

Sobre o diagrama em si, cada classe possui três partes: título, os atributos,
que representam as características da classe em si, e seus métodos, que são as
operações que o sistema pode operar nos objetos das classes \cite{guedes:2022}. Na
Figura~\ref{fig:classes}, pode-se observar os elementos descritos acima:

\begin{figure}[h]
\centering
\caption{Diagrama de Classes}
\label{fig:classes}
\includegraphics[width=0.9\linewidth]{classes-v2.png}
\source{Elaborado pelo autor (2026)}
\end{figure}

Essa representação apresenta alguns vários elementos importantes. Primeiro, um
acoplamento entre \textit{Usuário} e as demais classes. Isso é natural, e
desejado, considerando que é o único ator. Existe também a separação entre as
entidades como Gatilho, Humor, muito embora exista uma relação de associação
entre elas, através da classe denominada ``VinculoGatilhoHumor'', visto que ela
seria representada por uma tabela auxiliar ligando as duas, aliadas do atributo
de estimativa de impacto.

Necessário destacar aqui que a classe ``AcaoCuidado'' é uma classe abstrata,
representando uma generalização. Essa classe serve como um guarda-chuva para
representar diversos tipos de ações que o usuário pode realizar. O papel dela
aqui se torna guardar elementos em comum entre esses tipos de ações, seus
respectivos vínculos com gatilhos ou humores. A partir dela então, traça-se
outras sub-classes que contem um melhor contexto sobre a atividade em si, por
exemplo, na classe ``Consulta'' guarda-se dados como a duração da consulta, que
tipo de consulta, e potenciais anotações ou comentários em cima da mesma.
Entretanto, uma classe ``RegistroMedicamento'' não precisa guardar mais do que
a relação com qual medicamento e quando ele foi de fato tomado.

Por fim, existe diversos tipos de enumeradores que delimitam exatamente quais
tipos de valores podem ser escolhidos. Esses enumeradores servem também como um
mecanismo de validação, sendo movidos tanto para o próprio banco de dados e
posteriormente sendo implementados como forma de controle em outros pontos como
na API.

Ressalta-se, contudo, que o diagrama não contempla, padrões arquiteturais
complementares, como \textit{Service Layer}, \textit{Repository} e \textit{Data
Access Object (DAO)}, os quais são incorporados na etapa de implementação para
garantir uma separação adequada de responsabilidades e facilitar a manutenção e
escalabilidade do código, conforme as propostas de \citet{Fowler_2013} e
\citet{Gamma:2011}. Esses elementos são importantes, considerando os princípios
de legibilidade, simplicidade e responsabilidades, muitos organizados por
\citet{Martin_2012}, asseguram que o código produzido seja limpo e sustentável.
